\documentclass[12pt,a4paper]{article}
\usepackage[utf8]{inputenc}
\usepackage[T2A]{fontenc}
\usepackage[russian]{babel}
\usepackage{graphicx}
\usepackage{booktabs}
\usepackage{longtable}
\usepackage{geometry}
\usepackage{float}
\usepackage{paracol}
\usepackage{src/title}
\usepackage[hidelinks]{hyperref}
\usepackage{listings}
\usepackage{xcolor}
\usepackage{amsmath}
\usepackage{tocloft}
\usepackage{tabularx}

\geometry{top=2cm, left=2cm, right=2cm, bottom=2cm}
\graphicspath{{../lab_3/graphics/}{../lab_4/graphics/}{../lab_5/graphics/}{../lab_6/graphics/}}

\lstset{
  language=Python,
  basicstyle=\footnotesize\ttfamily,
  keywordstyle=\color{blue},
  stringstyle=\color{red},
  commentstyle=\color{gray}\itshape,
  tabsize=2,
  columns=fullflexible,
  keepspaces=true,
  frame=single,
}

\title{Модуль 2. Общий отчёт по лабораторным работам №3--№6: методы обучения с учителем}
\author{Ануфриев А.С. 465029}
\supervisor{Болдырева Е. А.}
\date{08.11.2026}

\abstract{
В отчёте обобщены результаты четырёх лабораторных работ модуля 2 по линейной регрессии, методу k-ближайших соседей (k-NN), деревьям решений и логистической регрессии. Все модели реализованы ``с нуля'' на языке Python с использованием только NumPy и Pandas без готовых библиотек машинного обучения. Для каждой работы приведено описание метода, псевдокод, результаты экспериментов на реальных датасетах (жильё Калифорнии, вино, грибы UCI, пассажиры Титаника), графики и выводы. В заключении дан сравнительный анализ методов: их особенности, преимущества и ограничения. В приложении приводятся фрагменты исходного кода реализации.
}
\keywords{линейная регрессия, k-NN, дерево решений, логистическая регрессия, NumPy, Python, классификация, регрессия}

\begin{document}
\maketitle
\tableofcontents
\newpage

\section{Введение}
Модуль 2 посвящён основным методам (supervised) обучения c учителем: регрессии и классификации. В ходе четырёх лабораторных работ (лабораторные №3--№6) были реализованы и исследованы четыре метода:

\begin{enumerate}
\item линейная регрессия методом наименьших квадратов (лабораторная №3);
\item метод k-ближайших соседей k-NN (лабораторная №4);
\item дерево решений (лабораторная №5);
\item логистическая регрессия (лабораторная №6).
\end{enumerate}

Общая особенность всех работ --- реализация алгоритмов собственным кодом на Python с использованием только NumPy и Pandas. Готовые средства машинного обучения (sklearn и аналоги) не применялись. В каждой работе проведена загрузка и предобработка реального датасета, разделение на обучающую и тестовую выборки, обучение модели, оценка качества и оформление результатов в виде отдельного отчёта с графиками. Ниже по каждой лабораторной работе приведены цель, описание метода, псевдокод, результаты и выводы, а затем --- сравнительный анализ всех четырёх методов.

\section{Лабораторная работа 1. Линейная регрессия (данные о жилье Калифорнии)}

\subsection{Цели и задачи}
Цель работы --- построение и сравнение моделей линейной регрессии для прогнозирования стоимости жилья в Калифорнии методом наименьших квадратов без применения готовых библиотек регрессии. Основные задачи:
\begin{enumerate}
\item получить и визуализировать описательную статистику по датасету (количество, среднее, стандартное отклонение, минимум, максимум, квантили);
\item выполнить предобработку данных: обработку отсутствующих значений, кодирование категориальных признаков, нормировку;
\item разделить данные на обучающую и тестовую выборки;
\item реализовать линейную регрессию методом наименьших квадратов собственным кодом (минимизация суммы квадратов разностей фактических и предсказанных значений), используя только NumPy и Pandas;
\item построить три модели с разными наборами признаков;
\item оценить каждую модель коэффициентом детерминации $R^2$ и сравнить результаты;
\item в качестве бонуса ввести синтетический признак.
\end{enumerate}

\subsection{Описание метода}
Линейная множественная регрессия описывает зависимость результативного признака $y$ от нескольких факторов $x_1,\dots,x_p$:
\begin{equation}
\hat{y} = a + b_1 x_1 + \dots + b_p x_p ,
\end{equation}
где $a$ --- свободный член, $b_j$ --- регрессионные коэффициенты. Метод наименьших квадратов (МНК) определяет оценки параметров вектором $B = (a,b_1,\dots,b_p)^\top$ как
\begin{equation}
B = (X^\top X)^{-1} X^\top Y ,
\end{equation}
где $X$ --- матрица признаков с первым столбцом единиц (свободный член), $Y$ --- вектор наблюдений целевой переменной. Качество модели оценивается коэффициентом детерминации
\begin{equation}
R^2 = 1 - \frac{\sum_i (y_i - \hat{y}_i)^2}{\sum_i (y_i - \bar{y})^2} .
\end{equation}

\subsection{Псевдокод метода}
\begin{lstlisting}
# X: feature matrix, y: target vector
# add column of ones for the intercept term
X_aug = column_stack([ones(len(X)), X])
# OLS coefficients: B = (X^T X)^{-1} X^T y
B = inv(X_aug.T @ X_aug) @ X_aug.T @ y
# prediction
y_pred = X_aug_test @ B
# quality: coefficient of determination
R2 = 1 - sum((y_true - y_pred)^2) / sum((y_true - mean(y_true))^2)
\end{lstlisting}

\subsection{Исходные данные}
Исходный датасет содержит 17~000 записей о жилых районах Калифорнии. Все 9 признаков числовые, пропущенных значений и категориальных признаков в исходных данных отсутствуют. Целевая переменная --- \emph{median\_house\_value} (медианная стоимость дома, тыс. долл.). Особенность данных: значения стоимости ограничены сверху отметкой 500~000 (цenzурное отсечение); таких записей 22 (0,13 \%), доход у них разный, поэтому это не выброс, подлежащий удалению, а техническая особенность датасета.

\subsection{Ход работы}
\textbf{Описательная статистика.} Вычислены размер выборки, среднее, стандартное отклонение, минимум, максимум и квантили 25\%, 50\%, 75\% для всех признаков. Для целевой переменной построена гистограмма распределения, для пары ``доход населения --- стоимость жилья'' --- диаграмма рассеяния (рисунки \ref{fig:lab3_hist}, \ref{fig:lab3_scatter}).

\begin{table}[H]
\centering
\caption{Описательная статистика по датасету о жилье Калифорнии.}
\label{tab:lab3_stats}
\small
\begin{tabular}{lrrrrrrrr}
\toprule
Признак & count & mean & std & min & 25\% & 50\% & 75\% & max \\
\midrule
longitude & 17000 & -119.56 & 2.01 & -124.35 & -121.79 & -118.49 & -118.00 & -114.31 \\
latitude & 17000 & 35.63 & 2.14 & 32.54 & 33.93 & 34.25 & 37.72 & 41.95 \\
housing\_median\_age & 17000 & 28.59 & 12.59 & 1.0 & 18.0 & 29.0 & 37.0 & 52.0 \\
total\_rooms & 17000 & 2643.66 & 2179.95 & 2.0 & 1462.0 & 2127.0 & 3151.25 & 37937.0 \\
total\_bedrooms & 17000 & 539.41 & 421.50 & 1.0 & 297.0 & 434.0 & 648.25 & 6445.0 \\
population & 17000 & 1429.57 & 1147.85 & 3.0 & 790.0 & 1167.0 & 1721.0 & 35682.0 \\
households & 17000 & 501.22 & 384.52 & 1.0 & 282.0 & 409.0 & 605.25 & 6082.0 \\
median\_income & 17000 & 3.88 & 1.91 & 0.50 & 2.57 & 3.54 & 4.77 & 15.00 \\
median\_house\_value & 17000 & 207300.91 & 115983.76 & 14999.0 & 119400.0 & 180400.0 & 265000.0 & 500001.0 \\
\bottomrule
\end{tabular}
\end{table}

\begin{figure}[H]
\centering
\includegraphics[width=0.55\linewidth]{target_hist.pdf}
\caption{Гистограмма распределения медианной стоимости жилья (лаб.~3).}
\label{fig:lab3_hist}
\end{figure}

\begin{figure}[H]
\centering
\includegraphics[width=0.55\linewidth]{price_vs_income.pdf}
\caption{Диаграмма рассеяния стоимости жилья от дохода населения (лаб.~3).}
\label{fig:lab3_scatter}
\end{figure}

\textbf{Предобработка.} Пропущенных значений не обнаружено, категориальных признаков нет. Все числовые признаки нормированы стандартизацией (вычет среднего, деление на стандартное отклонение); целевая переменная не нормировалась.

\textbf{Разделение на выборки.} Данные перемешаны и разделены на обучающую (13~600 записей, 80\%) и тестовую (3~400 записей, 20\%) с фиксированным зерном генератора random\_state~=~42.

\textbf{Реализация.} Коэффициенты регрессии вычислены собственным кодом на Python с использованием NumPy по формуле $B=(X^\top X)^{-1}X^\top Y$, что эквивалентно минимизации суммы квадратов разностей фактических и предсказанных значений. Готовые средства регрессии не использовались.

\textbf{Три модели.} Построены три модели с разными наборами признаков. В третью модель включён синтетический признак $rooms\_per\_household = total\_rooms / households$ (бонусное задание).

\begin{table}[H]
\centering
\caption{Сравнение трёх моделей линейной регрессии (признаки нормированы).}
\label{tab:lab3_models}
\small
\begin{tabular}{lll}
\toprule
Модель & Признаки & $R^2$ на тестовой выборке \\
\midrule
Модель 1 & median\_income & 0.4978 \\
Модель 2 & median\_income, housing\_median\_age & 0.5376 \\
Модель 3 & median\_income, total\_rooms, households, rooms\_per\_household & 0.5227 \\
\bottomrule
\end{tabular}
\end{table}

\begin{table}[H]
\centering
\caption{Оценочные коэффициенты моделей (включая свободный член, признаки нормированы).}
\label{tab:lab3_coef}
\small
\begin{tabular}{lrrr}
\toprule
 & Модель 1 & Модель 2 & Модель 3 \\
\midrule
Свободный член & 207066.73 & 206962.89 & 207033.39 \\
median\_income & 79512.76 & 81929.71 & 89613.01 \\
housing\_median\_age & --- & 21519.53 & --- \\
total\_rooms & --- & --- & -54374.09 \\
households & --- & --- & 56290.68 \\
rooms\_per\_household & --- & --- & -6.13 \\
\bottomrule
\end{tabular}
\end{table}

\begin{figure}[H]
\centering
\includegraphics[width=0.32\linewidth]{pred_1.pdf}
\includegraphics[width=0.32\linewidth]{pred_2.pdf}
\includegraphics[width=0.32\linewidth]{pred_3.pdf}
\caption{Фактические против предсказанных значений на тестовой выборке: Модель 1, 2, 3 (лаб.~3).}
\label{fig:lab3_pred}
\end{figure}

\subsection{Выводы}
Наилучшее качество показала Модель 2 (доход населения и возраст жилого дома), $R^2 \approx 0.54$ на тестовой выборке. Доход населения --- сильнейший отдельный признак ($R^2 \approx 0.50$ в одномерной модели). Добавление возраста дома повысило качество, а в третью модель количество комнат, число домохозяйств и синтетический признак не улучшили результат ($R^2 \approx 0.52$) из-за сильной линейной зависимости (мультиколлинеарности) между этими признаками.

\section{Лабораторная работа 2. Метод k-ближайших соседей (датасет о вине)}

\subsection{Цели и задачи}
Цель работы --- исследование и применение метода k-ближайших соседей (k-NN) для классификации датасета о вине. Основные задачи:
\begin{enumerate}
\item описательная статистика и визуализация признаков (включая 3D-визуализацию);
\item предобработка данных: пропуски, кодирование категориальных признаков, масштабирование;
\item реализация k-NN собственным кодом без готовых библиотек классификации;
\item постройка двух моделей: со случайным и с фиксированным наборами признаков;
\item оценка качества на тестовой выборке при различных $k$ и построение матриц ошибок.
\end{enumerate}

\subsection{Описание метода}
Метод k-ближайших соседей относит объект к классу, наиболее часто встречающемуся среди $k$ ближайших соседей в пространстве признаков. Расстояние вычисляется по правилу Евклида:
\begin{equation}
d(\mathbf{x}_i, \mathbf{x}_j) = \sqrt{\sum_{p=1}^{P} (x_{ip} - x_{jp})^2}.
\end{equation}
Класс определяется majoritaire голосованием среди $k$ ближайших обучающих примеров. Качество оценивается долей верных ответов (accuracy) и матрицей ошибок. k-NN --- ``ленивый'' алгоритм: этапа обучения в классическом смысле нет, все вычисления выполняются на этапе предсказания.

\subsection{Псевдокод метода}
\begin{lstlisting}
# X_train, y_train: training set; X: test objects
# pairwise Euclidean distances: ||a-b||^2 = ||a||^2 + ||b||^2 - 2 a*b
d = sqrt(sum(X^2, axis=1) + sum(X_train^2, axis=1).T - 2 * X @ X_train.T)
# k nearest neighbours for each test object
nearest = argsort(d, axis=1)[:, :k]
# class: majority vote among neighbour labels
y_pred[i] = argmax(bincount(y_train[nearest[i]]))
\end{lstlisting}

\subsection{Исходные данные}
Исходный датасет содержит 178 записей о вине. Признаков: 13 числовых (Alcohol, Malic Acid, Ash, Alcalinity of ash, Magnesium, Total phenols, Flavanoids, Nonflavanoid phenols, Proanthocyanins, Color intensity, Hue, OD280/OD315 of diluted wines, Proline). Целевая переменная --- Wine (три класса: 1, 2, 3). Пропущенных значений и категориальных признаков не обнаружено.

\subsection{Ход работы}
\textbf{Предобработка.} Все 13 числовых признаков подвергнуты стандартизации. Данные разделены на обучающую (142 записи, 80\%) и тестовую (36 записей, 20\%) с random\_state = 42.

\textbf{Реализация.} k-NN реализован на NumPy: парные евклидовые расстояния вычисляются матрично, затем сортировка и выбор $k$ ближайших соседей, класс определяется голосованием. Готовые средства классификации не использовались.

\textbf{Модели.} Построены две модели:
\begin{itemize}
\item Модель 1 --- признаки отбираются случайно (5 из 13: Alcohol, Color intensity, OD280/OD315 of diluted wines, Proanthocyanins, Total phenols);
\item Модель 2 --- фиксированный набор признаков (Alcohol, Malic Acid, Ash, Alcalinity of ash, Magnesium, Total phenols).
\end{itemize}
Оценка при $k = 3, 5, 10, 15$.

\begin{table}[H]
\centering
\caption{Точность k-NN на тестовой выборке.}
\label{tab:lab4_metrics}
\small
\begin{tabular}{llcccc}
\toprule
Модель & Признаки & k=3 & k=5 & k=10 & k=15 \\
\midrule
Модель 1 & случайные (5 признаков) & 0.9444 & 0.9444 & 0.9167 & 0.9167 \\
Модель 2 & фиксированные (6 признаков) & 0.8889 & 0.8889 & 0.9444 & 0.9444 \\
\bottomrule
\end{tabular}
\end{table}

\begin{table}[H]
\centering
\caption{Матрица ошибок, Модель 1, k=5 (строки --- истинные классы, столбцы --- предсказанные).}
\label{tab:lab4_cm1}
\small
\begin{tabular}{lccc}
\toprule
 & \multicolumn{3}{c}{Предсказанный класс} \\
Истинный класс & 1 & 2 & 3 \\
\midrule
1 & 14 & 0 & 0 \\
2 & 1 & 12 & 1 \\
3 & 0 & 0 & 8 \\
\bottomrule
\end{tabular}
\end{table}

\begin{table}[H]
\centering
\caption{Матрица ошибок, Модель 2, k=10.}
\label{tab:lab4_cm2}
\small
\begin{tabular}{lccc}
\toprule
 & \multicolumn{3}{c}{Предсказанный класс} \\
Истинный класс & 1 & 2 & 3 \\
\midrule
1 & 13 & 0 & 1 \\
2 & 0 & 14 & 0 \\
3 & 0 & 1 & 7 \\
\bottomrule
\end{tabular}
\end{table}

\begin{figure}[H]
\centering
\includegraphics[width=0.45\linewidth]{confusion_model1.pdf}
\includegraphics[width=0.45\linewidth]{confusion_model2.pdf}
\caption{Матрицы ошибок: Модель 1 (слева, k=5) и Модель 2 (справа, k=10) (лаб.~4).}
\label{fig:lab4_cm}
\end{figure}

\subsection{Выводы}
Обе модели обеспечивают высокую точность на тестовой выборке (до 0.9444 из 36 объектов). Модель 1 стабильна при $k=3$ и $k=5$ (0.9444), Модель 2 достигает наилучшего результата при $k=10$ и $k=15$ (0.9444). Основные ошибки приходятся на путаницу между классами 2 и 3, которые ближе друг к другу в пространстве признаков.

\section{Лабораторная работа 3. Деревья решений (датасет грибов UCI Mushroom)}

\subsection{Цели и задачи}
Целью работы является реализация алгоритма построения дерева решений для задачи бинарной классификации грибов (съедобный / ядовитый) и оценка его качества. Основные этапы:
\begin{enumerate}
\item загрузка и предобработка исходных данных;
\item случайный отбор $\sqrt{n}$ признаков;
\item построение небинарного дерева решений собственным кодом без сторонних библиотек для построения дерева;
\item оценка по Accuracy, precision и recall без использования библиотек;
\item построение кривых AUC-ROC и AUC-PR без использования библиотек.
\end{enumerate}

\subsection{Описание метода}
Дерево решений --- иерархическая структура правил классификации: каждый внутренний узел содержит проверку по какому-то признаку, ребра соответствуют возможным значениям признака, а листья содержат метки классов. Построение выполняется рекурсивно сверху вниз (подход алгоритма C4.5). В каждом узле выбирается признак с максимальным нормированным приростом информации:
\begin{equation}
\mathrm{GainRatio}(X) = \frac{\mathrm{Info}(T) - \mathrm{Info}_X(T)}{\mathrm{splitinfo}(X)},
\end{equation}
где $\mathrm{Info}(T) = -\sum_j p_j \log_2 p_j$ --- энтропия множества примеров, $\mathrm{Info}_X(T) = \sum_i \frac{|T_i|}{|T|}\mathrm{Info}(T_i)$ --- условная энтропия после разбиения по признаку $X$, $\mathrm{splitinfo}(X) = -\sum_i \frac{|T_i|}{|T|}\log_2\frac{|T_i|}{|T|}$. По выбранному признаку создаётся столько потомков, сколько у него значений (дерево не бинарное). Лист получает класс большинства объектов в узле.

\subsection{Псевдокод метода}
\begin{lstlisting}
build_tree(X, y, features, depth):
    node = majority class of y; share of poisonous = mean(y == 1)
    if all examples one class or depth = max or no features left:
        node.is_leaf = True; return node
    best feature = argmax information gain (GainRatio) over features
    if no gain: node.is_leaf = True; return node
    for each value v of best feature:
        child[v] = build_tree(X[feature==v], y[feature==v], features without best, depth+1)
    return node

predict(node, object):
    descend to leaf; return label or share of poisonous (probability)
\end{lstlisting}

\subsection{Исходные данные}
Используется датасет UCI Mushroom (\emph{agaricus-lepiota.data}, семейства \textit{Agaricus} и \textit{Lepiota}, источник --- справочник \emph{The Audubon Society Field Guide to North American Mushrooms}, 1981). Датасет содержит 8124 объекта, каждый описывается классом (съедобный \texttt{e} / ядовитый \texttt{p}) и 22 номинативными признаками; пропуски \emph{?} встречаются только в признаке \emph{stalk-root}. После удаления объектов с пропусками остаётся 5644 записи: 3488 съедобных (61.8\%) и 2156 ядовитых (38.2\%).

\subsection{Ход работы}
\begin{enumerate}
\item Чтение данных; исключение записей с пропусками в \emph{stalk-root}.
\item Случайный отбор $\sqrt{22} \approx 4$ признаков с зерном 42: \emph{cap-shape}, \emph{cap-surface}, \emph{gill-color}, \emph{stalk-color-above-ring}.
\item Стратифицированное разделение 70/30: обучение 3951 объект, тест 1693 объекта, соотношение классов сохранено.
\item Построение дерева на обучении: максимальная глубина 12, фактическая глубина 4, 85 внутренних узлов, 63 листа.
\item Метки классов и вероятности (доля ядовитых в листе) для теста.
\item Accuracy, precision, recall, AUC-ROC, AUC-PR собственным кодом (метод трапеций для площадей под кривыми).
\end{enumerate}

\begin{table}[H]
\centering
\caption{Матрица ошибок и метрики качества на тестовой выборке (1693 объекта, положительный класс --- ядовитый гриб).}
\label{tab:lab5_metrics}
\begin{tabular}{lrrr}
\toprule
 & Значение & & \\
\midrule
TP (яды распознаны) & 472 & FP (съедобные как яды) & 71 \\
TN (съедобные распознаны) & 975 & FN (яды как съедобные) & 175 \\
\midrule
Accuracy & \multicolumn{3}{r}{0.8547} \\
Precision & \multicolumn{3}{r}{0.8692} \\
Recall & \multicolumn{3}{r}{0.7295} \\
AUC-ROC & \multicolumn{3}{r}{0.9416} \\
AUC-PR & \multicolumn{3}{r}{0.9223} \\
\bottomrule
\end{tabular}
\end{table}

\begin{figure}[H]
\centering
\includegraphics[width=0.48\linewidth]{roc_plot.pdf}
\quad
\includegraphics[width=0.48\linewidth]{pr_plot.pdf}
\caption{ROC- и PR-кривые классификатора на тестовой выборке (лаб.~5): AUC-ROC = 0.9416, AUC-PR = 0.9223.}
\label{fig:lab5_curves}
\end{figure}

\subsection{Выводы}
Реализовано небинарное дерево решений без сторонних библиотек. На 4 случайно отобранных признаках дерево глубины 4 достигает Accuracy 0.8547, precision 0.8692 и recall 0.7295; AUC-ROC = 0.9416, AUC-PR = 0.9223. Получение вероятностей из долей объектов в листьях позволяет построить ROC- и PR-кривые. Ограничение: recall 0.7295 означает, что часть ядовитых грибов классифицируется как съедобные, что критично для задачи безопасного потребления.

\section{Лабораторная работа 4. Логистическая регрессия (датасет пассажиров Титаника)}

\subsection{Цели и задачи}
Целью работы является реализация модели логистической регрессии для бинарной классификации без использования сторонних библиотек машинного обучения, исследование влияния гиперпараметров на качество модели и оценка её работы на тестовых данных.
\begin{enumerate}
\item загрузка и предобработка датасета пассажиров Титаника;
\item описательная статистика и её визуализация;
\item разделение на обучающую и тестовую выборки;
\item реализация логистической регрессии ``с нуля'' (сигмоида, log-loss, градиентный спуск, настраиваемая скорость обучения и число итераций);
\item исследование гиперпараметров: скорость обучения, число итераций, метод оптимизации (градиентный спуск и оптимизация Ньютона);
\item оценка по accuracy, precision, recall, F1-Score на тестовой выборке.
\end{enumerate}

\subsection{Описание метода}
Логистическая регрессия моделирует вероятность принадлежности объекта к положительному классу с помощью сигмоидной функции:
\begin{equation}
\sigma(z) = \frac{1}{1 + e^{-z}}, \qquad p = \sigma(Xw),
\end{equation}
где $X$ --- матрица признаков со столбцом смещения, $w$ --- вектор весов, $p$ --- предсказанная вероятность. Порог классификации обычно равен 0.5: $p \ge 0.5 \Rightarrow$ класс 1.

Функция потерь (логарифмические потери, кросс-энтропия):
\begin{equation}
L(w) = -\frac{1}{N}\sum_{i=1}^{N}\left[y_i\ln p_i + (1-y_i)\ln(1-p_i)\right].
\end{equation}
Обучение --- градиентный спуск:
\begin{equation}
w \leftarrow w - \frac{\eta}{N}X^\top(p-y),
\end{equation}
где $\eta$ --- скорость обучения. Для оптимизации Ньютона используется гессиан $H = \frac{1}{N}X^\top W X$, $W = \mathrm{diag}(p_i(1-p_i))$, и обновление $w \leftarrow w - H^{-1}\nabla L(w)$. Функция выпукла, поэтому градиентные методы сходятся к глобальному минимуму.

\subsection{Псевдокод метода}
\begin{lstlisting}
sigmoid(z) = 1 / (1 + exp(-clip(z, -500, 500)))
log_loss(y, p) = -mean(y*log(p) + (1-y)*log(1-p))

gradient_descent(X, y, lr, n_iter):
    w = zeros
    repeat n_iter times:
        p = sigmoid(X @ w)
        grad = X.T @ (p - y) / N
        w = w - lr * grad

Newton(X, y, n_iter):
    repeat n_iter times:
        p = sigmoid(X @ w)
        grad = X.T @ (p - y) / N
        W = diag(p * (1 - p))
        H = X.T @ W @ X / N
        w = w - solve(H, grad)

predict(X, w): return (sigmoid(X @ w) >= 0.5)
\end{lstlisting}

\subsection{Исходные данные и предобработка}
Используется датасет пассажиров Титаника, 891 запись. Целевой признак --- \emph{Survived} (0 --- не выжил, 1 --- выжил); доля выживших около 38,4\%. Числовые признаки: \emph{Pclass}, \emph{Age}, \emph{SibSp}, \emph{Parch}, \emph{Fare}; категориальные: \emph{Sex}, \emph{Embarked}. Признаки \emph{Name}, \emph{Ticket}, \emph{Cabin}, \emph{PassengerId} исключены. Пропуски: в \emph{Age} --- 177, в \emph{Embarked} --- 2, в \emph{Cabin} --- 687 (признак исключён).

Предобработка без утечки информации:
\begin{enumerate}
\item данные разделены на обучающую и тестовую 80/20 (712 / 179, зерно 42) \emph{до} любой предобработки;
\item статистики для заполнения пропусков и стандартизации вычислены только на обучающей выборке (медиана Age/Fare, мода Embarked, среднее и std числовых признаков);
\item те же статистики применены к тестовой выборке;
\item пропуски Age и Fare заполнены медианой обучения, Embarked --- модой обучения;
\item Sex закодирован (male = 0, female = 1), Embarked --- one-hot кодированием (порты S, C, Q);
\item числовые признаки стандартизированы статистиками обучения.
\end{enumerate}
После предобработки в модели используется 9 признаков (5 числовых + пол + 3 порта one-hot).

\begin{table}[H]
\centering
\caption{Описательная статистика по датасету Титаника (все 891 запись после заполнения пропусков).}
\label{tab:lab6_stats}
\begin{tabular}{lrrrrrrrr}
\toprule
Признак & Count & Mean & Std & Min & 25\% & 50\% & 75\% & Max \\
\midrule
Survived & 891 & 0.3838 & 0.4866 & 0.0 & 0.0 & 0.0 & 1.0 & 1.0 \\
Pclass & 891 & 2.3086 & 0.8361 & 1.0 & 2.0 & 3.0 & 3.0 & 3.0 \\
Age & 891 & 29.3616 & 13.0197 & 0.42 & 22.0 & 28.0 & 35.0 & 80.0 \\
SibSp & 891 & 0.5230 & 1.1027 & 0.0 & 0.0 & 0.0 & 1.0 & 8.0 \\
Parch & 891 & 0.3816 & 0.8061 & 0.0 & 0.0 & 0.0 & 0.0 & 6.0 \\
Fare & 891 & 32.2042 & 49.6934 & 0.0 & 7.91 & 14.45 & 31.0 & 512.33 \\
\bottomrule
\end{tabular}
\end{table}

\begin{figure}[H]
\centering
\includegraphics[width=0.6\textwidth]{lab6_survived.png}
\caption{Распределение целевой переменной Survived (лаб.~6).}
\label{fig:lab6_survived}
\end{figure}

\subsection{Исследование гиперпараметров}
Проведено исследование: три значения скорости обучения (0.01, 0.1, 0.5), три числа итераций (100, 500, 1000) и два метода оптимизации (градиентный спуск, оптимизация Ньютона).

\begin{table}[H]
\centering
\caption{Качество модели на тестовой выборке для комбинаций гиперпараметров.}
\label{tab:lab6_hp}
\small
\begin{tabular}{lrrrrrrr}
\toprule
Метод & lr & Iter. & Loss & Acc & Prec. & Rec. & F1 \\
\midrule
GD & 0.01 & 100 & 0.6340 & 0.6983 & 0.6341 & 0.4000 & 0.4906 \\
GD & 0.01 & 500 & 0.5490 & 0.7207 & 0.6744 & 0.4462 & 0.5370 \\
GD & 0.01 & 1000 & 0.5115 & 0.7430 & 0.6863 & 0.5385 & 0.6034 \\
GD & 0.1 & 100 & 0.5118 & 0.7430 & 0.6863 & 0.5385 & 0.6034 \\
GD & 0.1 & 500 & 0.4503 & 0.7765 & 0.6866 & 0.7077 & 0.6970 \\
GD & 0.1 & 1000 & 0.4449 & 0.7765 & 0.6866 & 0.7077 & 0.6970 \\
GD & 0.5 & 100 & 0.4503 & 0.7765 & 0.6866 & 0.7077 & 0.6970 \\
GD & 0.5 & 500 & 0.4442 & 0.7765 & 0.6866 & 0.7077 & 0.6970 \\
GD & 0.5 & 1000 & 0.4442 & 0.7765 & 0.6866 & 0.7077 & 0.6970 \\
Newton & --- & 100 & 0.4442 & 0.7765 & 0.6866 & 0.7077 & 0.6970 \\
Newton & --- & 500 & 0.4442 & 0.7765 & 0.6866 & 0.7077 & 0.6970 \\
Newton & --- & 1000 & 0.4442 & 0.7765 & 0.6866 & 0.7077 & 0.6970 \\
\bottomrule
\end{tabular}
\end{table}

Наилучший результат по F1-Score показывает градиентный спуск с learning rate 0.1 и 500 итерациями: accuracy 0.7765, precision 0.6866, recall 0.7077, F1-Score 0.6970 (TP = 46, TN = 93, FP = 21, FN = 19). Увеличение learning rate свыше 0.1 и числа итераций свыше 500 не улучшает метрики: loss стабилизируется. Метод Ньютона сходится за немногие итерации, но не даёт преимущества в точности на тесте по сравнению с градиентным спуском.

\begin{table}[H]
\centering
\caption{Матрица ошибок лучшей модели (GD, lr = 0.1, 500 итераций) на тестовой выборке.}
\label{tab:lab6_confusion}
\begin{tabular}{lrr}
\toprule
 & Предсказано: не выжил & Предсказано: выжил \\
\midrule
На самом деле: не выжил & TN = 93 & FP = 21 \\
На самом деле: выжил & FN = 19 & TP = 46 \\
\bottomrule
\end{tabular}
\end{table}

\begin{figure}[H]
\centering
\includegraphics[width=0.6\textwidth]{lab6_learning_curve.png}
\caption{Кривая обучения (градиентный спуск, lr = 0.1, 500 итераций) (лаб.~6).}
\label{fig:lab6_curve}
\end{figure}

\begin{figure}[H]
\centering
\includegraphics[width=0.7\textwidth]{lab6_hyperparam_heatmap.png}
\caption{Зависимость F1-Score от learning rate и числа итераций (градиентный спуск) (лаб.~6).}
\label{fig:lab6_heatmap}
\end{figure}

\subsection{Выводы}
Реализована логистическая регрессия ``с нуля'' на Python (NumPy/Pandas): сигмоида, log-loss, градиентный спуск с настраиваемой скоростью обучения и числом итераций, а также оптимизация Ньютона. Проведена предобработка без утечки данных. Категориальный признак Embarked zakодирован one-hot кодированием. Наилучшее качество по F1-Score (0.6970) достигается при learning rate 0.1, 500 итерациях и градиентном спуске. Классы несбалансированы (доля выживших 38,4\%), что снижает качество предсказания меньшего класса при общем высоком accuracy.

\section{Сравнительный анализ методов}

\subsection{Сравнительная таблица}
\begin{table}[H]
\centering
\caption{Сравнение четырёх методов модуля 2.}
\label{tab:compare}
\small
\begin{tabularx}{\linewidth}{>{\raggedright\arraybackslash}p{2.2cm}>{\raggedright\arraybackslash}X>{\raggedright\arraybackslash}p{3.6cm}>{\raggedright\arraybackslash}X>{\raggedright\arraybackslash}X}
\toprule
Метод & Тип задачи & Датасет, качество (тест) & Преимущества & Ограничения \\
\midrule
Линейная регрессия (МНК) & Регрессия (непрерывная цель) & Жильё Калифорнии, $R^2$ до 0.54 & Простота, интерпретируемость коэффициентов, аналитическое решение, быстрота & Линейная зависимость, чувствительность к мультиколлинеарности и выбросам \\
k-ближайших соседей & Классификация (и регрессия) & Вино, accuracy до 0.9444 & Простота, нет этапа обучения, непараметричность, работает с любым числом классов & Медленное предсказание, ``проклятие размерности'', чувствительность к масштабу признаков и выбору $k$ \\
Дерево решений & Классификация & Грибы UCI, accuracy 0.8547, AUC-ROC 0.9416 & Интерпретируемость, номинативные признаки, нелинейные границы, мало подготовки данных & Склонность к переобучению, нестабильность, жадный алгоритм \\
Логистическая регрессия & Бинарная классификация & Титаник, accuracy 0.7765, F1 0.6970 & Вероятностный вывод, быстрая сходимость, хорошая база & Линейная граница разделения, сложности с нелинейными зависимостями \\
\bottomrule
\end{tabularx}
\end{table}

\subsection{Обсуждение}
Все четыре метода показали приемлемое качество на своих задачах. Линейная регрессия и логистическая регрессия --- линейные модели: первые предсказывают непрерывное значение через среднеквадратичную функцию потерь, вторые через сигмоиду и кросс-энтропию выдают вероятность класса. k-NN и дерево решений --- непараметрические методы, capables capturing нелинейных зависимостей без явной функциональной формы.

k-NN не требует обучения, но стоимость предсказания растет с размером выборки; качество сильно зависит от масштабирования признаков (в работе без стандартизации признаки с большими значениями, например Proline ~747, доминировали бы в расстоянии). Дерево решений естественно работает с номинативными признаками и даёт легко интерпретируемые правила, но может переобучаться (в работе ограничение глубины и выбор $\sqrt{n}$ признаков служили регуляризацией). Логистическая регрессия даёт калиброванные вероятности и работает лучше всего при линейно разделимых классах; при сильной нелинейности или большом числе взаимодействий признаков уступает деревьям и k-NN.

\subsection{Примеры лучшего применения каждого метода}
\begin{itemize}
\item \textbf{Линейная регрессия:} прогнозирование числовых величин, когда связь приближённо линейна (стоимость недвижимости, спрос, продажи), когда нужна интерпретируемость влияния каждого фактора.
\item \textbf{k-NN:} задачи классификации на небольших и средних выборках с осмысленной метрикой расстояния, системы рекомендаций, задачи без четкой функциональной зависимости между признаками и целевой переменной.
\item \textbf{Деревья решений:} смешанные числовые и категориальные данные, задачи, важна интерпретируемость решений (кредитный скоринг, медицинская диагностика), когда не требуется строгая параметрическая форма зависимости.
\item \textbf{Логистическая регрессия:} бинарная классификация с вероятностным выводом (отток клиентов, скоринг рисков, выживание), особенно как базовая модель, которую легко обучать и интерпретировать.
\end{itemize}

\section{Заключение}
В рамках модуля 2 реализованы четыре метода надзоруемого обучения --- линейная регрессия, k-NN, дерево решений и логистическая регрессия --- полностью собственным кодом на Python (NumPy/Pandas) без готовых библиотек машинного обучения. По каждому методу получены результаты на реальных датасетах: $R^2 \approx 0.54$ для регрессии стоимости жилья, accuracy до 0.9444 при k-NN на данных о вине, accuracy 0.8547 с AUC-ROC 0.9416 при классификации грибов, accuracy 0.7765 и F1 0.6970 для логистической регрессии на датасете Титаника. Линейные модели просты и интерпретируемы, но предполагают линейную зависимость; непараметрические методы (k-NN, деревья) лучше улавливают нелинейные зависимости, но требуют тщательной предобработки и защиты от переобучения. Выбор метода зависит от природы задачи (регрессия против классификации), типа признаков, объема данных и требований к интерпретируемости.

\section{Приложение. Исходный код}
Ниже приводятся ключевые фрагменты реализации каждого метода (полные скрипты хранятся в папках \texttt{prog/} соответствующих лабораторных работ).

\subsection{Линейная регрессия (лаб.~3, \texttt{lab3\_regression.py})}
\begin{lstlisting}
def preprocess(df):
    # check missing values, encode categorical features,
    # standardize numeric features (target not normalized)
    report = {}
    report["missing_total"] = int(df.isna().sum().sum())
    categorical_cols = df.select_dtypes(exclude=[np.number]).columns.tolist()
    ...
    feature_cols = [c for c in df.columns if c != TARGET]
    means = df[feature_cols].mean()
    stds = df[feature_cols].std().replace(0, 1.0)
    df[feature_cols] = (df[feature_cols] - means) / stds
    return df, report

def ols_fit(X, y):
    # B = (X^T X)^{-1} X^T y, first column of ones added here
    X_aug = np.column_stack([np.ones(len(X)), X])
    XtX = X_aug.T @ X_aug
    XtX_inv = np.linalg.inv(XtX)
    B = XtX_inv @ X_aug.T @ y
    return B

def ols_predict(X, B):
    X_aug = np.column_stack([np.ones(len(X)), X])
    return X_aug @ B

def r2_score(y_true, y_pred):
    ss_res = float(np.sum((y_true - y_pred) ** 2))
    ss_tot = float(np.sum((y_true - np.mean(y_true)) ** 2))
    return 1.0 - ss_res / ss_tot if ss_tot != 0 else float("nan")
\end{lstlisting}

\subsection{k-ближайших соседей (лаб.~4, \texttt{knn\_wine.py})}
\begin{lstlisting}
def train_test_split(X, y, test_size=0.2, random_state=42):
    rng = np.random.RandomState(random_state)
    n = X.shape[0]
    indices = rng.permutation(n)
    n_test = int(round(n * test_size))
    return X[indices[n_test:]], X[indices[:n_test]], y[indices[n_test:]], y[indices[:n_test]]

class KNNClassifier:
    def __init__(self, k=3):
        self.k = k
    def fit(self, X, y):
        self.X_train = np.asarray(X, dtype=float)
        self.y_train = np.asarray(y)
    def predict(self, X):
        X = np.asarray(X, dtype=float)
        # ||a-b||^2 = ||a||^2 + ||b||^2 - 2 a*b
        a2 = np.sum(X ** 2, axis=1, keepdims=True)
        b2 = np.sum(self.X_train ** 2, axis=1, keepdims=True).T
        ab = X @ self.X_train.T
        dists = np.sqrt(np.maximum(a2 + b2 - 2 * ab, 0.0))
        k_nearest = np.argsort(dists, axis=1)[:, : self.k]
        k_labels = self.y_train[k_nearest]
        # majority vote among k neighbours
        predictions = np.array([
            np.bincount(row.astype(int)).argmax() for row in k_labels
        ])
        return predictions

def confusion_matrix(y_true, y_pred, classes):
    n = len(classes)
    cm = np.zeros((n, n), dtype=int)
    idx = {c: i for i, c in enumerate(classes)}
    for t, p in zip(y_true, y_pred):
        cm[idx[t], idx[p]] += 1
    return cm
\end{lstlisting}

\subsection{Дерево решений (лаб.~5, \texttt{lab5\_tree.py})}
\begin{lstlisting}
def entropy(y):
    if len(y) == 0:
        return 0.0
    p = np.mean(y == 1)
    if p == 0.0 or p == 1.0:
        return 0.0
    return -p * np.log2(p) - (1.0 - p) * np.log2(1 - p)

def information_gain(col, y, values):
    total = len(y)
    h_parent = entropy(y)
    weighted = 0.0
    for v in values:
        mask = col == v
        w = np.sum(mask)
        if w == 0:
            continue
        weighted += (w / total) * entropy(y[mask])
    return h_parent - weighted

def build_tree(X, y, feature_indices, depth):
    node = Node()
    node.label = majority_label(y)
    node.pos_frac = np.mean(y == 1)
    # stopping conditions: pure node, max depth, features exhausted, few examples
    if (len(np.unique(y)) == 1 or depth >= MAX_DEPTH
            or len(feature_indices) == 0 or len(y) < 2):
        node.is_leaf = True
        return node
    best_gain = -1.0
    best_feat = None
    for fi in feature_indices:
        gain = information_gain(X[:, fi], y, np.unique(X[:, fi]))
        if gain > best_gain:
            best_gain = gain
            best_feat = fi
    if best_feat is None or best_gain <= 0.0:
        node.is_leaf = True
        return node
    node.feature_idx = best_feat
    remaining = [f for f in feature_indices if f != best_feat]
    for v in np.unique(X[:, best_feat]):
        mask = X[:, best_feat] == v
        node.children[v] = build_tree(X[mask], y[mask], remaining, depth + 1)
    return node

def predict_proba(node, x):
    # share of poisonous in the leaf the object falls into
    cur = node
    while not cur.is_leaf:
        v = x[cur.feature_idx]
        if v in cur.children:
            cur = cur.children[v]
        else:
            break
    return cur.pos_frac

def auc_roc(y_true, score):
    order = np.argsort(-score, kind="mergesort")
    y_sorted = y_true[order]
    P = np.sum(y_true == 1)
    Nn = np.sum(y_true == 0)
    tp = fp = 0
    points = [(0.0, 0.0)]
    for i in range(len(y_sorted)):
        if y_sorted[i] == 1:
            tp += 1
        else:
            fp += 1
        points.append((fp / Nn, tp / P))
    points.append((1.0, 1.0))
    # area by trapezoidal rule
    uniq = {}
    for fpr, tpr in points:
        uniq[fpr] = max(uniq.get(fpr, 0.0), tpr)
    xs = sorted(uniq.keys())
    auc = 0.0
    for i in range(1, len(xs)):
        auc += (xs[i] - xs[i - 1]) * (uniq[xs[i]] + uniq[xs[i - 1]]) / 2.0
    return auc, points
\end{lstlisting}

\subsection{Логистическая регрессия (лаб.~6, \texttt{lab6\_logreg.py})}
\begin{lstlisting}
def sigmoid(z):
    return 1.0 / (1.0 + np.exp(-np.clip(z, -500, 500)))

def log_loss(y_true, p):
    eps = 1e-12
    p = np.clip(p, eps, 1.0 - eps)
    return -float(np.mean(y_true * np.log(p) + (1.0 - y_true) * np.log(1.0 - p)))

def gradient_descent(X, y, lr, n_iter):
    w = np.zeros((X.shape[1], 1))
    history = []
    for _ in range(n_iter):
        p = sigmoid(X @ w)
        grad = X.T @ (p - y) / X.shape[0]
        w = w - lr * grad
        history.append(log_loss(y, p))
    return w, history

def newton_method(X, y, n_iter):
    w = np.zeros((X.shape[1], 1))
    history = []
    for _ in range(n_iter):
        p = sigmoid(X @ w)
        grad = X.T @ (p - y) / X.shape[0]
        W = np.diag((p * (1.0 - p)).ravel())
        H = (X.T @ W @ X) / X.shape[0]
        step = np.linalg.solve(H, grad)
        w = w - step
        history.append(log_loss(y, p))
    return w, history

def predict(X, w):
    return (sigmoid(X @ w) >= 0.5).astype(int)

def classification_metrics(y_true, y_pred):
    tp = int(np.sum((y_pred == 1) & (y_true == 1)))
    tn = int(np.sum((y_pred == 0) & (y_true == 0)))
    fp = int(np.sum((y_pred == 1) & (y_true == 0)))
    fn = int(np.sum((y_pred == 0) & (y_true == 1)))
    acc = (tp + tn) / max(len(y_true), 1)
    prec = tp / max(tp + fp, 1)
    rec = tp / max(tp + fn, 1)
    f1 = 2.0 * prec * rec / max(prec + rec, 1e-12)
    return acc, prec, rec, f1, tp, tn, fp, fn
\end{lstlisting}

\end{document}